\documentclass{amsart}
% For dealing with references we use the comment environment
\usepackage{verbatim}
\newenvironment{reference}{\comment}{\endcomment}
%\newenvironment{reference}{}{}
\newenvironment{slogan}{\comment}{\endcomment}
\newenvironment{history}{\comment}{\endcomment}

% For commutative diagrams we use Xy-pic
\usepackage[all]{xy}

% We use 2cell for 2-commutative diagrams.
\xyoption{2cell}
\UseAllTwocells

% We use multicol for the list of chapters between chapters
\usepackage{multicol}

% This is generally recommended for better output
\usepackage{lmodern}
\usepackage[T1]{fontenc}

% For cross-file-references

% Package for hypertext links:
\usepackage{hyperref}

% For any local file, say "hello.tex" you want to link to please

% Theorem environments.
%
\theoremstyle{plain}
\newtheorem{theorem}[subsection]{Theorem}
\newtheorem{proposition}[subsection]{Proposition}
\newtheorem{lemma}[subsection]{Lemma}

\theoremstyle{definition}
\newtheorem{definition}[subsection]{Definition}
\newtheorem{example}[subsection]{Example}
\newtheorem{exercise}[subsection]{Exercise}
\newtheorem{situation}[subsection]{Situation}

\theoremstyle{remark}
\newtheorem{remark}[subsection]{Remark}
\newtheorem{remarks}[subsection]{Remarks}

\numberwithin{equation}{subsection}

% Macros
%
\def\lim{\mathop{\mathrm{lim}}\nolimits}
\def\colim{\mathop{\mathrm{colim}}\nolimits}
\def\Spec{\mathop{\mathrm{Spec}}}
\def\Hom{\mathop{\mathrm{Hom}}\nolimits}
\def\Ext{\mathop{\mathrm{Ext}}\nolimits}
\def\SheafHom{\mathop{\mathcal{H}\!\mathit{om}}\nolimits}
\def\SheafExt{\mathop{\mathcal{E}\!\mathit{xt}}\nolimits}
\def\Sch{\mathit{Sch}}
\def\Mor{\mathop{\mathrm{Mor}}\nolimits}
\def\Ob{\mathop{\mathrm{Ob}}\nolimits}
\def\Sh{\mathop{\mathit{Sh}}\nolimits}
\def\NL{\mathop{N\!L}\nolimits}
\def\CH{\mathop{\mathrm{CH}}\nolimits}
\def\proetale{{pro\text{-}\acute{e}tale}}
\def\etale{{\acute{e}tale}}
\def\QCoh{\mathit{QCoh}}
\def\Ker{\mathop{\mathrm{Ker}}}
\def\Im{\mathop{\mathrm{Im}}}
\def\Coker{\mathop{\mathrm{Coker}}}
\def\Coim{\mathop{\mathrm{Coim}}}

% Boxtimes
%
\DeclareMathSymbol{\boxtimes}{\mathbin}{AMSa}{"02}

%
% Macros for moduli stacks/spaces
%
\def\QCohstack{\mathcal{QC}\!\mathit{oh}}
\def\Cohstack{\mathcal{C}\!\mathit{oh}}
\def\Spacesstack{\mathcal{S}\!\mathit{paces}}
\def\Quotfunctor{\mathrm{Quot}}
\def\Hilbfunctor{\mathrm{Hilb}}
\def\Curvesstack{\mathcal{C}\!\mathit{urves}}
\def\Polarizedstack{\mathcal{P}\!\mathit{olarized}}
\def\Complexesstack{\mathcal{C}\!\mathit{omplexes}}
% \Pic is the operator that assigns to X its picard group, usage \Pic(X)
% \Picardstack_{X/B} denotes the Picard stack of X over B
% \Picardfunctor_{X/B} denotes the Picard functor of X over B
\def\Pic{\mathop{\mathrm{Pic}}\nolimits}
\def\Picardstack{\mathcal{P}\!\mathit{ic}}
\def\Picardfunctor{\mathrm{Pic}}
\def\Deformationcategory{\mathcal{D}\!\mathit{ef}}

\setlength{\emergencystretch}{3em}
\begin{document}
\title[Stacks Project, Tag 056Y]{468 --- Fibrewise nonzerodivisors extend to flat relative Cartier slices for flat finitely presented morphisms}
\maketitle
\noindent Copyright (C) 2005--2025 Johan de Jong. Stacks Project authors. Source: \href{https://stacks.math.columbia.edu/tag/056Y}{Tag 056Y}.

\noindent Editorial difficulty note (not author text). Difficulty H1: The local equation must remain a nonzerodivisor on an actual neighbourhood and the hypersurface quotient must remain flat. The proof descends the equation and morphism to a Noetherian model, checks the fibre condition through a faithfully flat local map, and brings a uniformly localized exact flat sequence back to the original base..

\noindent Editorial provenance note (not author text). History: excerpt from revision \texttt{a0444\allowbreak 6e57e\allowbreak c1fbc\allowbreak 252a8\allowbreak 71afc\allowbreak ec775\allowbreak 2fb28\allowbreak 07b14}, more-morphisms.tex, lines 6165--6261. Reference presentation was adapted; original-author context and core support blocks were retained or added. The mathematical wording is unchanged. Licensed under GNU FDL 1.2 or later; no invariant sections or cover texts. See ../licenses/COPYING. Context blocks reproduce author definitions and supporting results; they are not additional counted problems.

\begin{lemma}
\label{lemma-slice-given-element}
Let $f : X \to S$ be a morphism of schemes.
Let $x \in X$ be a point with image $s \in S$.
Let $h \in \mathfrak m_x \subset \mathcal{O}_{X, x}$.
Assume
\begin{enumerate}
\item $f$ is locally of finite presentation,
\item $f$ is flat at $x$, and
\item the image $\overline{h}$ of $h$ in
$\mathcal{O}_{X_s, x} = \mathcal{O}_{X, x}/\mathfrak m_s\mathcal{O}_{X, x}$
is a nonzerodivisor.
\end{enumerate}
Then there exists an affine open neighbourhood $U \subset X$ of $x$
such that $h$ comes from $h \in \Gamma(U, \mathcal{O}_U)$ and such
that $D = V(h)$ is an effective Cartier divisor in $U$ with $x \in D$ and
$D \to S$ flat and locally of finite presentation.
\end{lemma}

\begin{proof}
We are going to prove this by reducing to the Noetherian case.
By openness of flatness (see
Theorem \href{https://stacks.math.columbia.edu/tag/0399}{theorem openness flatness (Tag 0399)})
we may assume, after replacing $X$ by an
open neighbourhood of $x$, that $X \to S$ is flat.
We may also assume that $X$ and $S$ are affine.
After possible shrinking $X$ a bit we may assume that there exists
an $h \in \Gamma(X, \mathcal{O}_X)$ which maps to our given $h$.

\medskip\noindent
We may write $S = \Spec(A)$ and we may write $A = \colim_i A_i$
as a directed colimit of finite type $\mathbf{Z}$ algebras.
Then by
Algebra, Lemma \href{https://stacks.math.columbia.edu/tag/02JO}{lemma flat finite presentation limit flat (Tag 02JO)}
or
Limits, Lemmas \href{https://stacks.math.columbia.edu/tag/01ZM}{lemma descend finite presentation (Tag 01ZM)},
\href{https://stacks.math.columbia.edu/tag/01ZN}{lemma descend affine finite presentation (Tag 01ZN)}, and
\href{https://stacks.math.columbia.edu/tag/01ZM}{lemma descend finite presentation (Tag 01ZM)}
we can find a cartesian diagram
$$
\xymatrix{
X \ar[r] \ar[d]_f & X_0 \ar[d]^{f_0} \\
S \ar[r] & S_0
}
$$
with $f_0$ flat and of finite presentation, $X_0$ affine, and
$S_0$ affine and Noetherian. Let $x_0 \in X_0$, resp.\ $s_0 \in S_0$
be the image of $x$, resp.\ $s$. We may also assume there exists an element
$h_0 \in \Gamma(X_0, \mathcal{O}_{X_0})$ which restricts to $h$ on $X$.
(If you used the algebra reference above then this is clear; if you used
the references to the chapter on limits then this follows from
Limits, Lemma \href{https://stacks.math.columbia.edu/tag/01ZM}{lemma descend finite presentation (Tag 01ZM)}
by thinking of $h$ as a morphism $X \to \mathbf{A}^1_S$.)
Note that $\mathcal{O}_{X_s, x}$ is a localization of
$\mathcal{O}_{(X_0)_{s_0}, x_0} \otimes_{\kappa(s_0)} \kappa(s)$, so that
$\mathcal{O}_{(X_0)_{s_0}, x_0} \to \mathcal{O}_{X_s, x}$ is a flat
local ring map, in particular faithfully flat. Hence the image
$\overline{h}_0 \in \mathcal{O}_{(X_0)_{s_0}, x_0}$
is contained in $\mathfrak m_{(X_0)_{s_0}, x_0}$ and is a nonzerodivisor.
We claim that after replacing $X_0$ by a principal open neighbourhood of
$x_0$ the element $h_0$ is a nonzerodivisor in
$B_0 = \Gamma(X_0, \mathcal{O}_{X_0})$ such that $B_0/h_0B_0$ is flat
over $A_0 = \Gamma(S_0, \mathcal{O}_{S_0})$.
If so then
$$
0 \to B_0 \xrightarrow{h_0} B_0 \to B_0/h_0B_0 \to 0
$$
is a short exact sequence of flat $A_0$-modules. Hence this remains exact
on tensoring with $A$ (by
Algebra, Lemma \href{https://stacks.math.columbia.edu/tag/00HL}{lemma flat tor zero (Tag 00HL)})
and the lemma follows.

\medskip\noindent
It remains to prove the claim above. The corresponding algebra statement
is the following (we drop the subscript ${}_0$ here):
Let $A \to B$ be a flat, finite type ring map of Noetherian rings.
Let $\mathfrak q \subset B$ be a prime lying over $\mathfrak p \subset A$.
Assume $h \in \mathfrak q$ maps to a nonzerodivisor in
$B_{\mathfrak q}/\mathfrak p B_{\mathfrak q}$.
Goal: show that after possible replacing $B$ by $B_g$ for some
$g \in B$, $g \not \in \mathfrak q$ the element $h$ becomes a nonzerodivisor
and $B/hB$ becomes flat over $A$. By
Algebra, Lemma \href{https://stacks.math.columbia.edu/tag/00MF}{lemma grothendieck (Tag 00MF)}
we see that $h$ is a nonzerodivisor in $B_{\mathfrak q}$ and that
$B_{\mathfrak q}/hB_{\mathfrak q}$ is flat over $A$.
By openness of flatness, see
Algebra, Theorem \href{https://stacks.math.columbia.edu/tag/00RC}{theorem openness flatness (Tag 00RC)}
or
Theorem \href{https://stacks.math.columbia.edu/tag/0399}{theorem openness flatness (Tag 0399)}
we see that $B/hB$ is flat over $A$ after replacing $B$ by $B_g$ for some
$g \in B$, $g \not \in \mathfrak q$. Finally, let
$I = \{b \in B \mid hb = 0\}$ be the annihilator of $h$. Then
$IB_{\mathfrak q} = 0$ as $h$ is a nonzerodivisor in $B_{\mathfrak q}$.
Also $I$ is finitely generated as $B$ is Noetherian.
Hence there exists a $g \in B$, $g \not \in \mathfrak q$ such that $IB_g = 0$.
After replacing $B$ by $B_g$ we see that $h$ is a nonzerodivisor.
\end{proof}
\end{document}
